% ------------------------------------------------------------------------------
% Metodologia
% ------------------------------------------------------------------------------

\chapter{Metodologia}
\label{chap_metodologia}
Cada capítulo deve conter uma pequena introdução (tipicamente, um ou dois parágrafos), que deve deixar claro o objetivo e o que será discutido no capítulo, bem como a organização do mesmo.

Este capítulo descreve os procedimentos adotados para a modelagem, simulação computacional e análise dos resultados do fenômeno da queda livre. São detalhadas as etapas de implementação do modelo matemático, a geração e organização dos dados, bem como os métodos utilizados para análise e validação dos resultados.

\section{Coleta e tratamento de dados}
\label{sec_coleta_e_tratamento_de_dados}

Os dados utilizados neste trabalho foram gerados por meio de simulações computacionais do modelo de queda livre, implementado em Python. Foram definidos parâmetros iniciais como altura, velocidade e aceleração da gravidade. O método de Euler foi empregado para discretizar as equações do movimento, permitindo o cálculo iterativo da posição e velocidade ao longo do tempo.

Todos os dados de entrada e saída das simulações foram salvos em arquivos \texttt{.txt} na pasta \texttt{src/dados}, garantindo reprodutibilidade e organização. Os gráficos gerados foram salvos em formato \texttt{.png} na pasta \texttt{src/graficos}, facilitando a análise visual dos resultados.

\section{Análise estatística}
\label{sec_analise_estatistica}

A análise dos resultados envolveu a comparação entre as soluções numéricas obtidas pelo método de Euler e a solução analítica do problema da queda livre. Foram avaliados diferentes valores de passo de tempo ($\Delta t$) para investigar o impacto na precisão dos resultados. O erro absoluto entre o tempo de queda simulado e o valor analítico foi calculado e representado graficamente, permitindo avaliar a acurácia do método numérico.

Além disso, foram realizadas simulações considerando a resistência do ar, possibilitando a comparação entre diferentes cenários físicos e a análise do comportamento do sistema sob diferentes hipóteses.
